\section{Code implementation}
In order to run the simulations for the three models and compute all the statistics presented above, a custom Python project was developed. It is structured into six modules:

\begin{itemize}
    \item \texttt{config.py}: consists of data classes storing the constant parameters of the simulation;
    \item \texttt{simulation.py}: contains the core logic of the simulation;
    \item \texttt{animation.py}, fully dedicated to the graphic representation of the flock dynamic;
    \item \texttt{statistics.py}: handles the computation of correlation and polarization functions;
    \item\texttt{auxiliary\_functions.py}: contains miscellaneous functions, mainly related to plotting;
    \item \texttt{main.py}: manages the options and the overall project workflow.
\end{itemize}

This section focuses exclusively on the \texttt{simulation.py} module, as it is the most relevant to the purpose of this report. For the sake of clarity, rather than performing a line-by-line code review, we will outline the fundamental logic of this module. \
At the beginning of the file, two classes are defined in order to handle the initialization phase:

\begin{itemize}
    \item \texttt{FlockState}: initializes the boids' positions and velocities;
    \item \texttt{PredatorState}: initializes the predator's position and velocity.
\end{itemize}

Following the initialization structures, the core computation function is presented:

\begin{minted}[breaklines, bgcolor=lightgray!20]{python}
def compute_distances_and_fov(pos, vel):
    dist_vects = pos[:, np.newaxis, :] - pos[np.newaxis, :, :]
    dist_norms = np.linalg.norm(dist_vects, axis=-1)
    dot_prods = (vel[:, np.newaxis, :] * (-dist_vects)).sum(axis=-1)
    vel_norms = np.linalg.norm(vel, axis=-1)[:, np.newaxis]
    cos_angles = dot_prods / np.maximum(vel_norms * dist_norms, 1e-9)
    return dist_vects, dist_norms, cos_angles
\end{minted}

This function computes the distance vectors, the distance norms, and the viewing angles between each pair of boids in the flock. Instead of employing a classical \texttt{for} loop to iterate over the elements of the flock's position matrix, the computation is highly optimized through \texttt{NumPy} array broadcasting. The objects returned by this function are extensively used throughout the code. \
Subsequently, three distinct functions are defined to simulate the tested models: \texttt{compute\_reynolds}, \texttt{compute\_couzin}, and \texttt{compute\_vicsek}. The architecture of these functions is structurally similar. First, a boolean mask matrix is defined to identify each boid's neighbors within its specific action range and field of view:

\begin{minted}[breaklines, bgcolor=lightgray!20]{python}
def compute_reynolds(pos, vel, dist_vects, dist_norms, cos_angles, predator_state):
    mask = (dist_norms < cfg.reynolds_const.action_range) & 
    (dist_norms > 0) & (cos_angles > cfg.glob_const.cos_fov)
    mask_3d = mask[:, :, np.newaxis]
    n_neighbors = mask.sum(axis=1)
    has_neighbors = n_neighbors > 0
    ...
\end{minted}

Then, the forces governing each model's interactions are computed according to their original mathematical formulations:

\begin{minted}[breaklines, bgcolor=lightgray!20]{python}
def compute_reynolds(pos, vel, dist_vects, dist_norms, cos_angles, predator_state):
    ...
    # Cohesion
    sum_pos = (pos[np.newaxis, :, :] * mask_3d).sum(axis=1)
    centroid = sum_pos[has_neighbors] / n_neighbors[has_neighbors, np.newaxis]
    neighb_dist = centroid - pos[has_neighbors]
    neighb_dist_norm = np.linalg.norm((neighb_dist), keepdims=True, axis=1)
    coh_vel[has_neighbors] = (neighb_dist) * \
        cfg.reynolds_const.coh_par / (neighb_dist_norm**3)
    # Alignement
    sum_vel = (vel[np.newaxis, :, :] * mask_3d).sum(axis=1)
    mean_vel = sum_vel[has_neighbors] / n_neighbors[has_neighbors, np.newaxis]
    ali_vel[has_neighbors] = (
        mean_vel - vel[has_neighbors]) * cfg.reynolds_const.ali_par
    # Separation
    safe_distance_sq = np.where(dist_norms == 0, 1.0, dist_norms**3)
    repulsion = dist_vects / safe_distance_sq[:, :, np.newaxis]
    sep_vel = (repulsion * mask_3d).sum(axis=1) * cfg.reynolds_const.sep_par
    ...
\end{minted}

Interactions related to obstacle and predator avoidance are handled globally by two separate functions, which can then be integrated into any of the three models:

\begin{minted}[breaklines, bgcolor=lightgray!20]{python}
def compute_obstacle_avoidance(flock_pos):
    obs_dist_vects = flock_pos[:, np.newaxis, :] -  cfg.obstacles_const.positions
    [np.newaxis, :, :]
    obs_dist_norms = np.linalg.norm(obs_dist_vects, axis=-1)
    mask = (obs_dist_norms < cfg.obstacles_const.action_range) & (
        obs_dist_norms > 0)
    mask_3d = mask[:, :, np.newaxis]
    safe_distance_sq = np.where(obs_dist_norms == 0, 1.0, obs_dist_norms**2)
    repulsion = obs_dist_vects / safe_distance_sq[:, :, np.newaxis]
    obs_avoid_vel = (repulsion * mask_3d).sum(axis=1) * \
        cfg.obstacles_const.rep_par
    return obs_avoid_vel
\end{minted}

\begin{minted}[breaklines, bgcolor=lightgray!20]{python}
def compute_predator_avoidance(flock_pos, pred_pos):
    pred_dist_vects = flock_pos[:, np.newaxis, :] - pred_pos[np.newaxis, :, :]
    pred_dist_norms = np.linalg.norm(pred_dist_vects, axis=-1)
    mask = (pred_dist_norms < cfg.predator_const.dist_par) & (
        pred_dist_norms > 0)
    mask_3d = mask[:, :, np.newaxis]
    safe_distance_sq = np.where(pred_dist_norms == 0, 1.0, pred_dist_norms**2)
    repulsion = pred_dist_vects / safe_distance_sq[:, :, np.newaxis]
    pred_avoid_vel = (repulsion * mask_3d).sum(axis=1) * \
        cfg.predator_const.sep_par
    return pred_avoid_vel
\end{minted}

Finally, each model processes and applies these spatial forces differently:

\begin{itemize}
    \item in \texttt{compute\_reynolds}, each acceleration request is accumulated and constrained within a finite capacity using a clamping function;
    \item in \texttt{compute\_couzin}, the zonal priority model is applied, and a spherically wrapped Gaussian noise is added to the target direction;
    \item in \texttt{compute\_vicsek}, the alignment force is applied to each boid, followed by the addition of uniform angular noise.
\end{itemize}

The final section of the \texttt{simulation.py} module is devoted to the \texttt{update\_flock} function. This function acts as the main controller, responsible for applying the selected method to compute the velocity updates between two subsequent time steps:

\begin{minted}[breaklines, bgcolor=lightgray!20]{python}
def update_flock(flock_state: FlockState, predator_state: PredatorState, method: str):
    dist_vects, dist_norms, cos_angles = compute_distances_and_fov(
        flock_state.pos, flock_state.vel)
    match method.lower():
        case "reynolds":
            boids_prov_vel = compute_reynolds(
                flock_state.pos, flock_state.vel, dist_vects, dist_norms, cos_angles, predator_state)
        case "couzin":
            boids_prov_vel = compute_couzin(
                flock_state.pos, flock_state.vel, dist_vects, dist_norms, cos_angles, predator_state)
        case "vicsek":
            boids_prov_vel = compute_vicsek(
                flock_state.pos, flock_state.vel, dist_norms, predator_state)
        case _:
            raise ValueError(
                f"Method '{method}' is invalid. Choose between 'reynolds', 'couzin' or 'vicsek'.")
    flock_state.vel += boids_prov_vel
    ...
\end{minted}

After the specific model's logic is computed, the boids' velocity vectors are updated and, if required by the model (e.g., Reynolds), clamped to enforce physical speed limits. Finally, an \texttt{if} statement is employed to manage the predator dynamics, updating its position based on the flock's centroid attraction while enforcing its own maximum speed constraints:

\begin{minted}[breaklines, bgcolor=lightgray!20]{python}
    ...
    if cfg.commands.predator_bool == True:
        pred_prov_vel = predator_move(
            flock_state.pos, predator_state.pos)
        predator_state.vel += pred_prov_vel
        # Predator's velocity limit
        pred_speed = np.linalg.norm(predator_state.vel, axis=1, keepdims=True)
        predator_state.vel = np.where(
            pred_speed > cfg.predator_const.max_speed,
            (predator_state.vel / pred_speed) * cfg.predator_const.max_speed,
            predator_state.vel
        )
        predator_state.pos += predator_state.vel

    flock_state.pos += flock_state.vel
\end{minted}
